\begin{example}
\label{example-affine-open-not-standard}
Consider the ring
$$
R = \{ f \in \mathbf{Q}[z]\text{ with }f(0) = f(1) \}.
$$
Consider the map
$$
\varphi : \mathbf{Q}[A, B] \to R
$$
defined by $\varphi(A) = z^2-z$ and $\varphi(B) = z^3-z^2$.  It is
easily checked that $(A^3 - B^2 + AB) \subset \Ker(\varphi)$ and that
$A^3 - B^2 + AB$ is irreducible. Assume that $\varphi$ is surjective;
then since $R$ is an integral domain (it is a subring of an integral
domain), $\Ker(\varphi)$ must be a prime ideal of $\mathbf{Q}[A, B]$.
The prime ideals which contain $(A^3-B^2 + AB)$ are $(A^3-B^2 + AB)$
itself and any maximal ideal $(f, g)$ with $f, g\in\mathbf{Q}[A, B]$
such that $f$ is irreducible mod $g$. But $R$ is not a field, so the
kernel must be $(A^3-B^2 + AB)$; hence $\varphi$ gives an isomorphism
$R \to \mathbf{Q}[A, B]/(A^3-B^2 + AB)$.

\medskip\noindent
To see that $\varphi$ is surjective, we must express any
$f\in R$ as a $\mathbf{Q}$-coefficient polynomial in $A(z) = z^2-z$
and $B(z) = z^3-z^2$. Note the relation $zA(z) = B(z)$. Let
$a = f(0) = f(1)$. Then $z(z-1)$ must divide $f(z)-a$, so we can write
$f(z) = z(z-1)g(z)+a = A(z)g(z)+a$.  If $\deg(g) < 2$, then
$h(z) = c_1z + c_0$ and $f(z) = A(z)(c_1z + c_0)+a = c_1B(z)+c_0A(z)+a$, so we
are done.  If $\deg(g)\geq 2$, then by the polynomial division
algorithm, we can write $g(z) = A(z)h(z)+b_1z + b_0$
($\deg(h)\leq\deg(g)-2$), so $f(z) = A(z)^2h(z)+b_1B(z)+b_0A(z)$.
Applying division to $h(z)$ and iterating, we obtain an expression
for $f(z)$ as a polynomial in $A(z)$ and $B(z)$; hence $\varphi$ is
surjective.

\medskip\noindent
Now let $a \in \mathbf{Q}$, $a \neq 0, \frac{1}{2}, 1$ and
consider
$$
R_a = \{ f \in \mathbf{Q}[z, \frac{1}{z-a}]\text{ with }f(0) = f(1)
\}.
$$
This is a finitely generated $\mathbf{Q}$-algebra as well: it is
easy to check that the functions $z^2-z$, $z^3-z$, and
$\frac{a^2-a}{z-a}+z$ generate $R_a$ as an $\mathbf{Q}$-algebra.  We
have the following inclusions:
$$
R\subset R_a\subset\mathbf{Q}[z, \frac{1}{z-a}], \quad
R\subset\mathbf{Q}[z]\subset\mathbf{Q}[z, \frac{1}{z-a}].
$$
Recall (Lemma \ref{lemma-spec-localization}) that for a ring T and a
multiplicative subset $S\subset T$, the ring map $T \to S^{-1}T$
induces a map on spectra $\Spec(S^{-1}T) \to \Spec(T)$
which is a homeomorphism onto the subset
$$
\{\mathfrak p \in \Spec(T) \mid S \cap \mathfrak p = \emptyset\}
\subset \Spec(T).
$$
When $S = \{ 1, f, f^2, \ldots\}$ for some $f\in T$, this is
the open set $D(f)\subset T$.  We now verify a corresponding
property for the ring map $R \to R_a$: we will show that the map
$\theta : \Spec(R_a) \to \Spec(R)$ induced by inclusion
$R\subset R_a$ is a homeomorphism onto an open subset of
$\Spec(R)$ by verifying that $\theta$ is an injective local
homeomorphism.  We do so with respect to an open cover of
$\Spec(R_a)$ by two distinguished opens, as we now describe.
For any $r\in\mathbf{Q}$, let $\text{ev}_r : R \to \mathbf{Q}$ be the
homomorphism given by evaluation at $r$.  Note that for $r = 0$ and
$r = 1-a$, this can be extended to a homomorphism
$\text{ev}_r' : R_a \to \mathbf{Q}$ (the latter because $\frac{1}{z-a}$
is well-defined at $z = 1-a$, since $a\neq\frac{1}{2}$).  However,
$\text{ev}_a$ does not extend to $R_a$.  Write
$\mathfrak{m}_r = \Ker(\text{ev}_r)$. We have
$$
\mathfrak{m}_0 = (z^2-z, z^3-z),
$$
$$
\mathfrak{m}_a = ((z-1 + a)(z-a), (z^2-1 + a)(z-a)), \text{ and}
$$
$$
\mathfrak{m}_{1-a} = ((z-1 + a)(z-a), (z-1 + a)(z^2-a)).
$$
To verify this, note that the right-hand sides are clearly contained in
the left-hand sides. Then check that the right-hand sides are
maximal ideals by writing the generators in terms of $A$ and $B$,
and viewing $R$ as $\mathbf{Q}[A, B]/(A^3-B^2 + AB)$. Note that
$\mathfrak{m}_a$ is not in the image of $\theta$: we have
$$
(z^2 - z)^2(z - a)\left(\frac{a^2 - a}{z - a} + z\right) =
(z^2 - z)^2(a^2 - a) + (z^2 - z)^2(z - a)z
$$
The left hand side is in $\mathfrak m_a R_a$ because
$(z^2 - z)(z - a)$ is in $\mathfrak m_a$ and because
$(z^2 - z)(\frac{a^2 - a}{z - a} + z)$ is in $R_a$. Similarly
the element $(z^2 - z)^2(z - a)z$ is in $\mathfrak m_a R_a$
because $(z^2 - z)$ is in $R_a$ and $(z^2 - z)(z - a)$ is in $\mathfrak m_a$.
As $a \not \in \{0, 1\}$ we conclude that
$(z^2 - z)^2 \in \mathfrak m_a R_a$. Hence
no ideal $I$ of $R_a$ can satisfy $I \cap R = \mathfrak m_a$, as such
an $I$ would have to contain $(z^2 - z)^2$, which is in $R$ but not in
$\mathfrak m_a$. The distinguished open set
$D((z-1 + a)(z-a))\subset\Spec(R)$ is equal to the complement of
the closed set $\{\mathfrak{m}_a, \mathfrak{m}_{1-a}\}$.
Then check that $R_{(z-1 + a)(z-a)} = (R_a)_{(z-1 + a)(z-a)}$; calling
this localized ring $R'$, then, it follows that the map $R \to R'$
factors as $R \to R_a \to R'$.  By Lemma
\ref{lemma-spec-localization}, then, these maps express
$\Spec(R') \subset \Spec(R_a)$ and
$\Spec(R') \subset \Spec(R)$ as open subsets; hence
$\theta : \Spec(R_a) \to \Spec(R)$, when restricted to
$D((z-1 + a)(z-a))$, is a homeomorphism onto an open subset.
Similarly, $\theta$ restricted to
$D((z^2 + z + 2a-2)(z-a)) \subset \Spec(R_a)$ is a homeomorphism
onto the open subset $D((z^2 + z + 2a-2)(z-a)) \subset \Spec(R)$.
Depending on whether $z^2 + z + 2a-2$ is irreducible or not over
$\mathbf{Q}$, this former distinguished open set has complement
equal to one or two closed points along with the closed point
$\mathfrak{m}_a$. Furthermore, the ideal in $R_a$ generated by the
elements $(z^2 + z + 2a-a)(z-a)$ and $(z-1 + a)(z-a)$ is all of $R_a$, so
these two distinguished open sets cover $\Spec(R_a)$. Hence in
order to show that $\theta$ is a homeomorphism onto
$\Spec(R)-\{\mathfrak{m}_a\}$, it suffices to show
that these one or two points can never equal $\mathfrak{m}_{1-a}$.
And this is indeed the case, since $1-a$ is a root of $z^2 + z + 2a-2$
if and only if $a = 0$ or $a = 1$, both of which do not occur.

\medskip\noindent
Despite this homeomorphism which mimics the behavior of a
localization at an element of $R$, while
$\mathbf{Q}[z, \frac{1}{z-a}]$ is the localization of $\mathbf{Q}[z]$
at the maximal ideal $(z-a)$, the ring $R_a$ is {\it not} a
localization of $R$: Any localization $S^{-1}R$ results in more
units than the original ring $R$.  The units of $R$ are
$\mathbf{Q}^\times$, the units of $\mathbf{Q}$.  In fact, it is
easy to see that the units of $R_a$ are $\mathbf{Q}^*$.
Namely, the units of $\mathbf{Q}[z, \frac{1}{z - a}]$ are
$c (z - a)^n$ for $c \in \mathbf{Q}^*$ and $n \in \mathbf{Z}$
and it is clear that these are in $R_a$ only if $n = 0$.
Hence $R_a$ has no more units than
$R$ does, and thus cannot be a localization of $R$.

\medskip\noindent
We used the fact that $a\neq 0, 1$ to ensure that
$\frac{1}{z-a}$ makes sense at $z = 0, 1$.  We used the fact that
$a\neq 1/2$ in a few places: (1) In order to be able to talk about
the kernel of $\text{ev}_{1-a}$ on $R_a$, which ensures that
$\mathfrak{m}_{1-a}$ is a point of $R_a$ (i.e., that $R_a$ is
missing just one point of $R$). (2) At the end in order to conclude
that $(z-a)^{k + \ell}$ can only be in $R$ for $k = \ell = 0$; indeed, if
$a = 1/2$, then this is in $R$ as long as $k + \ell$ is even. Hence
there would indeed be more units in $R_a$ than in $R$, and $R_a$
could possibly be a localization of $R$.
\end{example}